\textbf{а)}
Группа корней $n$ степени из $1$ изоморфна $\mathbb{Z}_n$.

Тогда примитивные корни $=$  элементы порядка $n$. Их, как мы знаем, $\varphi(n)$.

\textbf{б)}

\begin{itemize}
    \item $n = 2: x^2 - 1$. Корни: $1, -1$. Примитивные корни: $-1$, минимальный многочлен: $x+1$ 
    \item $n = 3: x^3 - 1$. Корни: $1, \omega, \omega^2$. Примитивные корни: $\omega, \omega^2$, минимальный многочлен: $x^2 + x + 1$
    \item $n = 4: x^4 - 1$. Корни: $1, -1, i, -i$. Примитивные корни: $i, -i$, минимальный многочлен: $x^2+1$
    \item $n = 5: x^5 - 1$. Корни: $1, \xi_5, \xi_5^2, \xi_5^3, \xi_5^4$. Примитивные корни: $\xi_5, \xi_5^2, \xi_5^3, \xi_5^4$, минимальный многочлен: $x^4 + x^3 + x^2 + 1$
    \item $n = 6: x^6 - 1$. Корни: $1, \omega, \omega^2, -1, -\omega, -\omega^2$. Примитивные корни: $-\omega, -\omega^2$, минимальный многочлен: $x^2 - x + 2$
\end{itemize}

\textbf{в)}

Пусть $n = p$, тогда примитивные корни -- все, кроме $1$. Поэтому минимальный многочлен это $\frac{x^p - 1}{x - 1} = x^{p-1} + \ldots + 1$